% 人工对照底本逐节点转写；独立TikZ重绘。

\begin{tikzpicture}[x=1cm,y=-1cm,every node/.style={anchor=west,align=left,inner sep=2pt,font=\fontsize{10}{13.0}\selectfont},edge/.style={draw=black!65,line width=.45pt}]
\node[text width=2.500cm,draw=black!65] (n1) at (0.000,1.329) {非法使用\\（盗用、冒用）};
\node[text width=4.500cm,draw=black!65] (n2) at (2.920,0.524) {银行卡（包括实体卡和信息资料）};
\node[text width=2.580cm,draw=black!65] (n3) at (7.840,0.524) {定信用卡诈骗罪};
\node[text width=2.932cm,draw=black!65] (n4) at (11.460,0.524) {例外：抢劫，盗窃};
\node[text width=3.283cm,draw=black!65] (n5) at (2.920,1.566) {他人的微信、支付宝};
\node[text width=1.526cm,draw=black!65] (n6) at (7.840,1.566) {定盗窃罪};
\node[text width=3.635cm,draw=black!65] (n7) at (2.920,2.371) {他人的蚂蚁花呗、借呗};
\node[text width=2.229cm,draw=black!65] (n8) at (7.840,2.371) {定贷款诈骗罪};
\draw[edge] (n3.east)--(11.310,0.524)--(11.310,0.524)--(n4.west);
\draw[edge] (n2.east)--(7.690,0.524)--(7.690,0.524)--(n3.west);
\draw[edge] (n1.east)--(2.770,1.329)--(2.770,0.524)--(n2.west);
\draw[edge] (n5.east)--(7.690,1.566)--(7.690,1.566)--(n6.west);
\draw[edge] (n1.east)--(2.770,1.329)--(2.770,1.566)--(n5.west);
\draw[edge] (n7.east)--(7.690,2.371)--(7.690,2.371)--(n8.west);
\draw[edge] (n1.east)--(2.770,1.329)--(2.770,2.371)--(n7.west);
\end{tikzpicture}

